\lecture{27}{Wed 16 Nov 2022 10:48}{}

\begin{remark}
	The situation $r = 0$ is important, where $g$ ``divides'' $f$, and $f(x) = \in g(x) F[x] = (g)$. $g$ divides $f$ iff $f$ lies in ideal generated by $g$.
\end{remark}
\begin{theorem}(Polynomials rings are PID)
	Suppose that $I\neq 0$ in an ideal of $F[x]$. Then  $I = (g) = \{f(x)g(x): f(x) \in F[x]\} $ for some $g \in F[x]$.
\end{theorem}
\begin{proof}
	Try to identiy a polynomial $g$ to generate $I$, by considering 
	\[
		d = \min \{\operatorname{deg} h: h \in I, \operatorname{deg} g \ge 0\} 
	.\] 
	and take $g$ to be any polynomial of degree $d$. Now, if $h \in I$ we have $\operatorname{deg} h \ge \operatorname{deg} g$, and by the division algorithm we have $h = q g + r$, where $q, r \in F[x]$ and $\operatorname{deg} r < \operatorname{deg} g$ or $r = 0$. But $h \in I$ and $q g \in I, so mkr = h - q g \in I$. 
	And so $I$ contains the polynomial $r$ which has degree smaller than $\operatorname{deg}  g$ (contradiction) or $r = 0$. Thus $h = q g \in (g)$.
\end{proof}

\begin{definition}[Principal Ideal Domain (PID)]
	A \textit{principal  ideal domain} is an integral domain $R$ satisfying the condition that whenever $I$ is an ideal of $R$, then $I = (a) = \{xa: x \in R\}$ for some $a \in I$.
\end{definition}

\begin{remark}
	Every principal ideal domain (PID) is a unique factorization domain (UFD)
\end{remark}
\begin{remark}
	Generator of an ideal need not be unique. But this can be resolved by focusing on \textit{monic} polynomials with lead coefficient $1$. To see this, note that if $I \triangleleft F[x]$ and $\left( p(x) \right)  = I = \left( m(x) \right) $, then since  $p(x) \in \left( m(x) \right) $, $p(x) = f(x) m(x)$, in particular  $\operatorname{deg}p \le  \operatorname{deg} m$. Similarly $\operatorname{deg} m \le \operatorname{deg}  p$. Hence $\operatorname{deg} p = \operatorname{deg}  m$, so $\operatorname{deg} f = 0$. Hence they are same up to scaling.
\end{remark}

\begin{definition}[Divisibility]
	\begin{enumerate}
		\item When $f, g, \in F[x]$ and $g \neq 0$, we say $g$ divides $f$ (and write $g | f$) if $f = ag $ for some $a \in F[x]$.
		\item The \textit{greatest common divisor} of $f $ and $g$ is the \textbf{monic} divisor $d \in F[x]$ of $f$ and $g$ divisible by all other such common divisors. This is written as $(f,g)$ or $\operatorname{gcd}(f,g)$.
	\end{enumerate}
\end{definition}
\begin{theorem}(Analog of consequence of Euclid's Algorithm)
	Let $f, g \in F[x]$ and suppose $g \neq 0$. Then $d = (f,g)$ exists and there exists $a, b \in F[x]$ such that $d = a f + b g$.
\end{theorem}
\begin{proof}
	Let $I = \{a f + b g: a, b \in F[x]\} $. We
	\begin{claim}
		$I \triangleleft F[x]$.
	\end{claim}
	To check this: whenever $a, b \in I$, say $a = u_1 f + v_1 g$ and $b = u_2 f + v_2 g$,  one has
	\[
		a + b = (u_1+u_2)f + (v_1+v_2) g \in I
	.\] 
	and for all $t \in F[x]$,
	\[
		ta = (tu_1)f + (tv_1)g \in I
	.\] 
	Also $I \neq 0$ since $f \in I$. Hence $I \triangleleft F[x]$.

	Then since $F[x]$ is a PID, we know that $I = (d)$ for some $d \in F[x]$. But then, since $f, g \in I$ we have $f $ and $g$ are both multiples of $d$. Then $d | f$ and $d | g$, so $d$ is a common divisor of $f$ and $g$.

	But the definition of $I$ shows that 
	\[
		d = u f + v g
	,\]
	for some $u,v \cap F[x]$. Thus, whenever $h | f$ and $h | g$, we have $h | d$. If we take (as we can) $d$ monic, then we have shown that $d = (f,g)$.
\end{proof}

From the proof we see that any generating monic polynomial is automatically GCD.

Say $f$ and $g$ are \textit{relatively prime} when $\left( f(x), g(x) \right)  = 1$, iff there exists  $u, v \in F[x]$ s.t. $uf + vg = 1$.

\begin{lemma}
	Suppose $\left( q(x), f(x) \right) =1$ and $q | fg$. Then $q | g$.
\end{lemma}
\begin{proof}
	There exists $a, b \in F[x]$ such that $a q + b f = 1$. Then 
	\[
		(aq)g + b(fg) = g
	.\] 
	Now the lhs is divisible by $q$ since $q | aq $ and $q | fg $. Hence $q | g$.
\end{proof}

\begin{lemma}
	Suppose that $\left( q(x), f(x) \right)  = 1$ and $q(x) | f(x) g(x)$. Then  $q(x) | g(x)$.
\end{lemma}
\begin{proof}
	We have that there exists $a, b \in F[x]$ such that $af + bq = 1$. Thus
	\[
		g = a(fg) + b(qg)
	\] 
	which is divisible by $q$.
\end{proof}


\begin{definition}
	We say $p \in F[x]$ is \textit{irreducible} if
	\begin{enumerate}
		\item $\operatorname{deg} p \ge 1$ 
		\item whenever $f \in F[x]$, either $p | f$ or $(p, f) = 1$.
	\end{enumerate}
\end{definition}
As a consequence, $p$ irreducible and $p = f g$ with $\operatorname{deg} f \ge \operatorname{deg} g\ge 1$, then either $f = c p$ or $g = c p$, some $c \in F$.
